\section{Computerization, Remote Work, and Competing Interpretations}\label{sec:competing}

The primary design conditions beta exposure on Webb's software-patent measure. Substituting other pre-existing technology measures changes the remaining occupational comparison. The beta coefficient is $-0.2085$ with O*NET computer-use importance, $-0.1512$ with O*NET computer-use level, $-0.1277$ with routine-task intensity, and $-0.1001$ with Frey--Osborne automation susceptibility, compared with $-0.1311$ under Webb. All corresponding intervals exclude zero, but the more than twofold coefficient range rules out the claim that ``controlling for computerization'' is a single invariant operation.

Remotability is strongly correlated with several exposure measures, especially AIOE, but correlation creates imprecision rather than automatic non-identification. Adding Dingel--Neiman occupational work-from-home feasibility leaves the beta estimate negative at $-0.1000$ with interval $[-0.1902,-0.0099]$. Under alpha, the estimate is $-0.0786$ with interval $[-0.1550,-0.0021]$. The remote coefficient changes sign across the two architectures. The narrow conclusion is that occupational remotability does not mechanically absorb the exposure tail contrast in these models. The exercise does not show that AI dominates remote work or that remote work is unimportant.

The categorical event study reports no detected differential pretrend under the existing common-multiplier max-statistic criterion: the largest absolute pre-period $t$ statistic is 1.502 and the wild-score $p$ value is .929. Forty of 42 observed post coefficients are negative, but the path is neither immediate nor monotone. A static post coefficient is an era average. Interest rates, technology-sector adjustment, return-to-office changes, post-pandemic normalization, and occupation-by-age shocks remain plausible concurrent forces.

The occupation-influence audit reinforces this limitation. Part of the extreme-tail contrast is supplied by favorable young-relative evolution in a low-exposure food-service occupation. This is consistent with the need to interpret the coefficient against broader post-pandemic labor-market normalization, but the deletion result does not identify normalization---or any other mechanism---as its cause.

These competing interpretations define rather than defeat the paper's contribution. The empirical object is a conditional association whose stability is examined across constructed exposure treatments. The analysis can establish that a sign persists across specified architectures, that magnitudes depend on conditioning choices, and that occupational rankings differ. It cannot convert exposure into adoption or separate AI causally from every contemporaneous shock.

